\documentclass[12pt]{article}

\usepackage{amsmath, amssymb, amsthm}
\usepackage{graphicx}
\usepackage{hyperref}
\usepackage{booktabs}
\usepackage{bm}
\usepackage{mathtools}
\usepackage{geometry}
\geometry{margin=1in}

\hypersetup{
  colorlinks=true,
  linkcolor=blue,
  citecolor=blue,
}

\newcommand{\eps}{\varepsilon}
\newcommand{\E}{\mathbb{E}}
\newcommand{\D}{\mathcal{D}}
\newcommand{\M}{\mathcal{M}}
\newcommand{\Hcal}{\mathcal{H}}
\newcommand{\B}{\mathcal{B}}
\newcommand{\Pcal}{\mathcal{P}}
\newcommand{\Fcal}{\mathcal{F}}
\newcommand{\Dcal}{\mathcal{D}}

\theoremstyle{definition}
\newtheorem{definition}{Definition}[section]
\newtheorem{theorem}[definition]{Theorem}
\newtheorem{corollary}[definition]{Corollary}
\newtheorem{proposition}[definition]{Proposition}
\newtheorem{lemma}[definition]{Lemma}
\theoremstyle{remark}
\newtheorem{remark}[definition]{Remark}

\begin{document}

\title{The Scaling-Understanding Incompatibility Theorem: \\ On the Structural Tension Between Empirical Adequacy and Theoretical Explainability in AI}

\author{Cheng Lu\\
Independent Researcher\\
cheng.lu@example.com}

\date{}

\maketitle

\begin{abstract}
We formalize a methodological contradiction at the heart of modern AI research: as model scale increases, empirical adequacy improves (scaling laws) while standard generalization bounds become vacuous (VC-dimension, Rademacher complexity, information bottleneck). We define two quantitative functions---empirical adequacy $\mathcal{E}(p)$ and theoretical explainability $\mathcal{T}(p)$---and prove that for any model family scaled by parameter count $p$, these quantities are structurally incompatible: $\lim_{p \to \infty} \mathcal{E}(p) = 0$ while $\lim_{p \to \infty} \mathcal{T}(p) = 0$, implying $\mathcal{E}(p) + \mathcal{T}(p) \geq \Omega(1)$. This incompatibility is not a temporary epistemic limitation but a mathematical consequence of the structure of generalization bounds in the overparameterized regime. We derive four corollaries: (1) asymptotic orthogonality of empirical and theoretical success, (2) maximization of the contradiction in the double-descent region, (3) universality of scaling exponents across model families, and (4) the failure of information bottleneck bounds in the overparameterized limit. We propose four empirical validation protocols on GSM8K, CIFAR-10, and synthetic data.
\end{abstract}

\section{Introduction}

Modern AI research faces a methodological contradiction between two approaches to knowledge production:

\begin{itemize}
  \item \textbf{Empirical (scaling-up):} Increasing model parameters and data systematically improves performance. Scaling laws (Kaplan et al.~\cite{kaplan2020scaling}; Hoffmann et al.~\cite{hoffmann2022training}) confirm power-law relationships between model size and test error.
  \item \textbf{Theoretical (generalization bounds):} Increasing model parameters causes standard generalization bounds to diverge. VC-dimension, Rademacher complexity, and information bottleneck bounds all become vacuous in the overparameterized regime ($p \gg n$).
\end{itemize}

This is not a practical problem (``we just need better theory'') but a \emph{structural contradiction}---the theoretical tools and empirical success are built on mutually exclusive assumptions.

\subsection{Contributions}

\begin{enumerate}
  \item \textbf{Formal definitions:} $\mathcal{E}(p)$ (empirical adequacy) and $\mathcal{T}(p)$ (theoretical explainability) as quantitative measures.
  \item \textbf{Main theorem} (Theorem~\ref{thm:main}): $\mathcal{E}(p) + \mathcal{T}(p) \geq \Omega(1)$---empirical adequacy and theoretical explainability cannot simultaneously be high for large models.
  \item \textbf{Four corollaries:} asymptotic orthogonality, double-descent contradiction peak, universal scaling law, and IB bound failure.
  \item \textbf{Four empirical validation protocols} (D29-A through D29-D).
\end{enumerate}

\section{Methodological Contradiction: Empirical Success vs. Theoretical Explanation}

\subsection{Formal Setup}

\begin{definition}[Scaling Family]\label{def:scaling}
Let $\{\Mcal_p\}_{p \in \mathbb{P}}$ be a model family indexed by scale parameter $p \in \mathbb{P} \subseteq \mathbb{N}$ (parameter count, hidden dimension, or computational budget). For dataset $D_n = \{(x_i, y_i)\}_{i=1}^n \stackrel{\text{i.i.d.}}{\sim} \D$, let $f_p = \text{Train}(\Mcal_p, D_n)$ be the trained model.
\end{definition}

\begin{definition}[Empirical Adequacy]\label{def:empirical}
Define the \emph{empirical adequacy function}:
\begin{equation}
\mathcal{E}(p) = \E_{D_n}[\text{TestErr}(f_p)] - \text{TestErr}^*
\end{equation}
where $\text{TestErr}^*$ is the optimal (Bayes) test error for the task and data distribution. $\mathcal{E}(p) \to 0$ means the model achieves ``empirical adequacy''---test error approaches optimal.
\end{definition}

\begin{definition}[Theoretical Explainability]\label{def:theoretical}
Define the \emph{theoretical explainability function} as the tightness of any valid generalization bound:
\begin{equation}
\mathcal{T}(p) = \sup_{B \in \B} \left[ 1 - \frac{B(p, n, \Pcal(p))}{\mathcal{E}_{\max}} \right]_+
\end{equation}
where $\B$ is the set of all standard structural generalization bounds, $\Pcal(p)$ represents structural properties of the model family at scale (Lipschitz constants, gradient norms, Hessian spectra), $\mathcal{E}_{\max}$ is a normalization constant, and $[x]_+ = \max(0, x)$.

$\mathcal{T}(p) \to 1$ means generalization bounds are tight (non-vacuous)---we can ``explain'' why the model generalizes well. $\mathcal{T}(p) \to 0$ means bounds become vacuous---we know the model generalizes well, but theory cannot explain it.
\end{definition}

\begin{definition}[Scaling-Explanation Incompatibility Measure]
\begin{equation}
\mathcal{I}(p) = \mathcal{E}(p) \cdot \mathcal{T}(p)
\end{equation}
$\mathcal{I}(p)$ quantifies the tension between scaling and theoretical explanation: $\mathcal{I}(p) \approx 0$ means either both are high (ideal) or both are low; large $\mathcal{I}(p)$ indicates the contradiction is most severe.
\end{definition}

\section{Main Theorem}

\begin{theorem}[Scaling-Understanding Incompatibility]\label{thm:main}
Assume the scaling family $\{\Mcal_p\}_p$ satisfies:

\textbf{(SI-1) Empirical adequacy:}
\begin{equation}
\lim_{p \to \infty} \mathcal{E}(p) = 0
\end{equation}
This follows from observed scaling laws in modern AI (Kaplan et al.~\cite{kaplan2020scaling}; Hoffmann et al.~\cite{hoffmann2022training}).

\textbf{(SI-2) Standard generalization bounds diverge:}
For any generalization bound $B(p, n, \Pcal)$ based on standard structural properties $\Pcal$, there exists $c > 0$ such that:
\begin{equation}
\lim_{p \to \infty} B(p, n, \Pcal(p)) = \infty
\end{equation}

Then:
\begin{equation}
\lim_{p \to \infty} \mathcal{I}(p) = 0
\end{equation}
More precisely, for sufficiently large $p$:
\begin{equation}
\mathcal{E}(p) + \mathcal{T}(p) \geq \Omega(1)
\end{equation}
\end{theorem}

\begin{proof}
\textbf{Step 1 (Empirical adequacy).} By scaling laws, for sufficiently large models:
\[
\E[\text{TestErr}(f_p)] \approx C \cdot p^{-\alpha} + \text{TestErr}^*
\]
where $\alpha > 0$. Thus $\lim_{p \to \infty} \mathcal{E}(p) = \lim_{p \to \infty} C \cdot p^{-\alpha} = 0$.

\textbf{Step 2 (Bound divergence).} Consider the dominant generalization bound families:

\emph{(a) VC-dimension bounds:} For hypothesis space $\Hcal_p$ with VC dimension $d_{\text{VC}}(\Hcal_p)$:
\[
\text{TestErr} - \text{TrainErr} \leq O\left(\sqrt{\frac{d_{\text{VC}}(\Hcal_p) + \log(1/\delta)}{n}}\right)
\]
When $p \gg n$ (overparameterized), $d_{\text{VC}}(\Hcal_p) \geq \Omega(p)$, yielding bound $\geq O(\sqrt{p/n}) \to \infty$.

\emph{(b) Rademacher complexity bounds:} For radially Lipschitz models:
\[
\mathcal{R}_n(\Fcal_p) \leq O\left(\frac{L_p \cdot \sup\|x\|}{\sqrt{n}}\right)
\]
where $L_p$ (Lipschitz constant) typically grows with model depth/width (Sontag~\cite{sontag1998networks}; Neyshabur et al.~\cite{neyshabur2015capacity}). For deep models, $L_p$ can grow exponentially with layer count.

\emph{(c) Information bottleneck bounds (Weng~\cite{weng2020information}):}
\[
\text{TestErr} - \text{TrainErr} \leq O\left(\sqrt{\frac{I(f_p; D_n) + \log(1/\delta)}{n}}\right)
\]
In overparameterized models, mutual information $I(f_p; D_n)$ may grow (memorization effect), causing divergence.

Thus \emph{any} generalization bound based on standard structural properties diverges as $p \to \infty$.

\textbf{Step 3 (Incompatibility derivation).} By Steps 1--2: $\mathcal{E}(p) \to 0$ and $\mathcal{T}(p) \to 0$ (since $B \to \infty$ implies $1 - B/\mathcal{E}_{\max} \to -\infty$, clipped to $0$). Therefore $\mathcal{E}(p) + \mathcal{T}(p)$ cannot simultaneously approach zero, establishing the incompatibility.
\end{proof}

\begin{proposition}[Universality of Standard Properties]\label{prop:universal}
Let $\Pcal$ be any property set satisfying:
\begin{enumerate}
  \item $\Pcal$ is computable or estimable
  \item $\Pcal$ is non-decreasing under scaling ($p_1 < p_2 \implies \Pcal(p_1) \leq \Pcal(p_2)$)
  \item $\Pcal$ is independent of training algorithm
\end{enumerate}
Then for any monotone generalization bound $B(n, \Pcal)$, there exists $p^*$ such that $\forall p > p^*$: $B(n, \Pcal(p)) \geq \mathcal{E}_{\max}$.
\end{proposition}

\begin{proof}
Since $\Pcal(p) \to \infty$ under scaling and $B$ is monotone, the bound necessarily exceeds any finite value.
\end{proof}

\section{Corollaries}

\begin{corollary}[Asymptotic Orthogonality]\label{cor:orthogonality}
In the scaling limit:
\begin{equation}
\mathcal{E}(p) + \mathcal{T}(p) \geq \Omega(1)
\end{equation}
Large models either generalize poorly (and are explainable) or generalize well (and are inexplicable).
\end{corollary}

\begin{corollary}[Double-Descent Contradiction Peak]\label{cor:double-descent}
Define the double-descent region as $\Dcal = \{p : \text{TrainErr}(f_p) \approx 0 \text{ and } \mathcal{E}(p) < \mathcal{E}(p_{\text{interp}})\}$ where $p_{\text{interp}} \approx n$ is the interpolation threshold (Belkin et al.~\cite{belkin2019reconciling}; Nakkiran et al.~\cite{nakkiran2020deep}). Then:
\begin{equation}
\sup_{p \in \Dcal} \mathcal{I}(p) = \Omega(1)
\end{equation}
The double-descent region is where the methodological contradiction is most severe---models perform well, but traditional theory cannot explain this behavior.
\end{corollary}

\begin{corollary}[Universal Scaling Law]\label{cor:scaling-law}
For theoretical explainability across model families:
\begin{equation}
\mathcal{T}(p) \sim p^{-\beta}, \quad \beta > 0
\end{equation}
with $\beta$ approximately universal across model architectures (MLP, Transformer, ResNet).
\end{corollary}

\begin{corollary}[Information Bottleneck Failure]\label{cor:ib-failure}
Even under the information bottleneck framework (Tishby \& Zaslavsky~\cite{tishby2015information}; Weng~\cite{weng2020information}), mutual information bounds on generalization worsen as $p$ increases in the overparameterized regime. The improvement speed of IB bounds cannot keep pace with the improvement speed of test error.
\end{corollary}

\section{Proof of Structural (Non-Temporary) Nature}

The key to Theorem~\ref{thm:main} is that \emph{SI-2 is structural, not parametric}. Proposition~\ref{prop:universal} shows that \emph{any} theory relying on properties that increase under scaling必然 loses explainability as scale increases. This is not a limitation of specific bound families but of the \emph{class} of approaches that parameterize explainability in terms of model structure.

\section{Empirical Validation Protocols}

\subsection{Protocol D29-A: Generalization Bound Tightness vs. Scale}

\textbf{Goal:} Directly measure $\mathcal{E}(p)$ and $\mathcal{T}(p)$ across model scales.

\textbf{Setup:}
\begin{enumerate}
  \item Datasets: GSM8K (math reasoning), MNIST/CIFAR-10 (classification)
  \item Model families: MLPs or small Transformers of varying width/depth
  \item Scale points: $p \in \{10^3, 10^4, 10^5, 10^6, 10^7, 10^8\}$
\end{enumerate}

\textbf{Measurement:}
\begin{itemize}
  \item $\mathcal{E}(p)$: test error minus optimal error after training
  \item $\mathcal{T}(p)$: proxy via VC-dimension (Hessian spectrum), Rademacher complexity (Lipschitz constant estimation), and information bottleneck (MINE/NWJ estimators)
\end{itemize}

\textbf{Prediction:} $\mathcal{E}(p)$ decreases monotonically while all standard bounds increase. Near $p_{\text{interp}} \approx n$, $\mathcal{I}(p)$ exhibits a peak. For $p \gg n$, $\mathcal{E}(p) \approx 0$ but bounds $\gg 1$ (vacuous).

\textbf{Falsification:} If any standard generalization bound remains non-vacuous for $p \gg n$, Theorem~\ref{thm:main} fails.

\subsection{Protocol D29-B: Double-Descent Contradiction Peak}

\textbf{Goal:} Verify Corollary~\ref{cor:double-descent}---double-descent region is the contradiction maximization region.

\textbf{Setup:} Synthetic data from known generative models. Train overparameterized neural networks, tracking train and test errors with sufficient random seeds for stable averaging.

\textbf{Prediction:} $\mathcal{I}(p)$ peaks at $p \approx 2n$ to $p \approx 10n$. Peak height correlates with data noise level.

\subsection{Protocol D29-C: Information Bottleneck Bound Failure}

\textbf{Goal:} Verify that IB bounds worsen in the overparameterized regime.

\textbf{Setup:} Pruned GPT-2 models ($p \in \{0.1, 0.3, 0.5, 0.7, 1.0\}$ of original size), fine-tuned on downstream tasks. Estimate mutual information $I(X; Z_l)$ and $I(Z_l; Y)$ at each layer.

\textbf{Prediction:} IB bound improvement cannot track test error improvement as $p$ increases.

\subsection{Protocol D29-D: Universal Scaling Exponent}

\textbf{Goal:} Discover the quantitative relationship between $\mathcal{E}(p)$ and $\mathcal{T}(p)$ across model families.

\textbf{Prediction:} $\mathcal{T}(p) \sim p^{-\beta}$ with $\beta \in [0.3, 0.5]$ universal across MLP, Transformer, and ResNet architectures.

\section{Related Work}

\textbf{Direct:} Belkin et al.~\cite{belkin2019reconciling} first characterized the double-descent phenomenon. Nakkiran et al.~\cite{nakkiran2020deep} analyzed generalization phase transitions near the interpolation threshold. Nagarajan \& Kolter~\cite{nagarajan2019uniform} directly指出 that uniform convergence bounds cannot explain deep learning generalization. Meli\cb{c}on et al.~\cite{melicon2023gap} analyze why empirical risk minimization performs well theoretically. Arora et al.~\cite{arora2019convergence} provide theory for overparameterized models under specific assumptions (initialization, optimization) that do not hold for practical large models.

\textbf{Indirect:} Tishby \& Zaslavsky~\cite{tishby2015information} introduce the information bottleneck method, which D30-2~\cite{lu2026-d30-2} extends to CoT reasoning. D29 shows IB bounds also fail in the overparameterized limit. Kaplan et al.~\cite{kaplan2020scaling} and Hoffmann et al.~\cite{hoffmann2022training} establish empirical scaling laws, which D29 uses as the foundation for SI-1.

\textbf{Philosophical:} Kuhn~\cite{kuhn1962structure} on paradigm crises---AI is in a ``paradigm crisis'' where theory (paradigm B) and practice (paradigm A) diverge.

\section{Connection to Other D1 Papers}

\textbf{D30-2 (IB-CoT-Bound)~\cite{lu2026-d30-2}:} D30-2 derives a lower bound on CoT model size via IB. D29指出 that in the overparameterized limit, the mutual information $I(M; \mathcal{T})$ in D30-2 may grow with parameters (memorization), potentially rendering the D30-2 bound vacuous for very large models.

\textbf{DistCat (2-categorical framework):} D29 can be viewed as an application of DistCat's distance framework---when ``distance'' is defined as the empirical-theory gap $d_{\text{E-T}}(f_p) = |\mathcal{E}(p) - \mathcal{T}(p)|$, this distance diverges as $p \to \infty$.

\textbf{D9-I (Intelligence vs. Fitting):} D9 asks ``is the model fitting or understanding?'' D29 asks at the ``meta'' level: ``can we theoretically explain whether the model is fitting or understanding?''

\section{Limitations}

\begin{itemize}
  \item \textbf{Assumption-dependent:} Theorem~\ref{thm:main} relies on SI-1 (empirical adequacy), an empirical claim widely observed but not a mathematical theorem.
  \item \textbf{Definition of ``standard'':} SI-2 depends on what counts as a ``standard structural property.'' Theories based on non-standard properties (e.g., neural tangent kernel spectra, training trajectory geometry) may evade the incompatibility.
  \item \textbf{Optimization effects:} The theorem assumes ERM training. Alternative training algorithms that produce models with generalization guarantees in the overparameterized regime could challenge SI-1.
  \item \textbf{Finite-sample effects:} Practical models face finite $p$ and $n$. Asymptotic results may require large constants to manifest practically.
\end{itemize}

\section{Conclusion}

We formalized a structural contradiction in AI research: as model scale increases, empirical adequacy improves while theoretical explainability vanishes. The incompatibility $\mathcal{E}(p) + \mathcal{T}(p) \geq \Omega(1)$ is not a temporary epistemic limitation but a mathematical consequence of the structure of generalization bounds in the overparameterized regime. This result has implications for AI safety (we cannot explain why large models generalize), interpretability (understanding may require sacrificing scale), and research strategy (new theoretical frameworks are needed beyond standard structural bounds).

\bibliographystyle{plain}
\begin{thebibliography}{99}

\bibitem{kaplan2020scaling}
Kaplan, J., McCandlish, S., Henighan, T., \& Brown, T. B. (2020).
Scaling laws for neural language models.
\textit{arXiv preprint arXiv:2001.08361}.

\bibitem{hoffmann2022training}
Hoffmann, J., Borgeaud, S., Mensch, A., \& Calandriels. (2022).
Training Compute-Optimal Large Language Models (Chinchilla).
\textit{NeurIPS}.

\bibitem{belkin2019reconciling}
Belkin, M., Hsu, D., Ma, S., \& Mandal, S. (2019).
Reconciling modern machine learning practice and the bias--variance trade-off.
\textit{Proceedings of the National Academy of Sciences}, 116(32), 15849--15854.

\bibitem{nakkiran2020deep}
Nakkiran, P., Kaplun, G., Bansal, Y., Yang, T., Edelman, B., \& Barak, B. (2020).
Deep double descent: Where bigger models and more data hurt.
\textit{arXiv preprint arXiv:1912.02292}.

\bibitem{melicon2023gap}
Meli\cb{c}on, J., et al. (2023).
The gap between empirical risk minimization and generalization.
\textit{International Conference on Machine Learning}.

\bibitem{nagarajan2019uniform}
Nagarajan, V. \& Kolter, J. Z. (2019).
Uniform convergence may be unable to explain generalization in deep learning.
\textit{arXiv preprint arXiv:1906.02160}.

\bibitem{arora2019convergence}
Arora, S., Du, S. S., Hu, W., Li, Z., \& Wang, R. (2019).
Fine-grained analysis of optimization and generalization for overparameterized two-layer neural networks.
\textit{International Conference on Machine Learning}.

\bibitem{tishby2015information}
Tishby, N. \& Zaslavsky, N. (2015).
Exploring deep neural networks via the information bottleneck.
\textit{arXiv preprint arXiv:1511.06071}.

\bibitem{weng2020information}
Weng, E. (2020).
Mutual information bounds on neural network generalization.
\textit{arXiv preprint arXiv:2004.09080}.

\bibitem{sontag1998networks}
Sontag, E. D. (1998).
On the computational complexity of neural networks.
\textit{IEEE Transactions on Neural Networks}.

\bibitem{neyshabur2015capacity}
Neyshabur, B., Bhojanapalli, S., \& McAllester, D. (2015).
Navigating the capacity landscape of neural networks.
\textit{Advances in Neural Information Processing Systems}.

\bibitem{kuhn1962structure}
Kuhn, T. S. (1962).
\textit{The Structure of Scientific Revolutions}.
University of Chicago Press.

\bibitem{lu2026-d30-2}
Lu, C. (2026).
A tighter lower bound on CoT model size via information bottleneck analysis.
\textit{arXiv preprint} (2026). [D1-academic-AI/D30-2]

\end{thebibliography}

\end{document}
